\begin{claim}{C.72}{Proved}{A spiral boundary obstruction with the full discrete local centralizer}
There is an explicit based loop of closed-disk homeomorphisms
fixing the boundary pointwise whose generated group $G$ is
transitive on the interior and has full interior centralizer
equal to a specified closed discrete free proper cyclic group
$Q=\langle T_*\rangle$. Every forward $Q$-orbit has the whole
boundary circle as its limit set. Nevertheless the actual
restriction map
\[
 C_{\mathcal D_\partial}(Q)\longrightarrow
              C_{\operatorname{Homeo}(\mathbb D)}(Q)
                                                               \tag{C.72.1}
\]
is not surjective: the proof displays an interior centralizer
map with no continuous extension at any boundary point.
Here $Q$ on the left means its specified closed-disk action.

The proof also supplies an explicit interior conjugacy
$j$ from $Q$ to a parabolic Fuchsian group. Under that
conjugacy the same obstructed map does extend as a
boundary-fixed homeomorphism. Thus the map
$h\mapsto jhj^{-1}$ between the two interior centralizers
does not identify their actual boundary-relative subgroups.
The exact failed inclusion is witnessed below; no
comparison of unnamed compactifications is used.

The original loop in this example has an explicit
boundary-fixed contraction. The result therefore
refutes the displayed extension comparison, including
the versions supplemented by transitivity and the
full-circle orbit limit set; it supplies no failure of
the level-zero loop comparison in C.12.
\end{claim}

\begin{proof}
\emph{The actual spiral coordinate and its inverse.}
Regard the real plane as $\mathbb C$, with
$z=x+iy$, and define
\[
 \begin{split}
 p(z)&=\frac{z}{1+|z|}
                \exp\bigl(i\log(1+|z|)\bigr),\\
 p^{-1}(w)&=\frac{w}{1-|w|}
                \exp\bigl(i\log(1-|w|)\bigr)
                            \qquad(w\in\mathbb D).
 \end{split}                                                   \tag{C.72.2}
\]
Both maps take zero to zero. The radial coordinates
are $r\mapsto r/(1+r)$ and
$\rho\mapsto\rho/(1-\rho)$, which are inverse.
The angular changes in the two displayed formulas
are opposite because $1+r=(1-\rho)^{-1}$.
They prove both inverse identities and continuity,
including at zero. Thus (C.72.2) is an actual
homeomorphism of the plane with the disk interior.
It is used to specify the action on the given
closed disk; no boundary extension of $p$ is asserted.

The following elementary estimate retains its boundary
topology. Suppose $|a|\leq M$ and $r=|z|>2M+1$.
Write $r'=|z+a|$. The reverse triangle inequality
gives $|r'-r|\leq M$, and
\[
 \begin{split}
 \left|\frac{z+a}{r'}-\frac z r\right|
                     &\leq\frac{2M}{r-M},\\
 |\log(1+r')-\log(1+r)|
                     &\leq\frac{M}{1+r-M},\\
 \left|\frac{r'}{1+r'}-\frac r{1+r}\right|
                     &\leq\frac{M}{(1+r-M)(1+r)} .
 \end{split}                                                   \tag{C.72.3}
\]
The first inequality follows by adding and subtracting
$z/r'$, and the second by the mean value theorem.
Combining them with $|e^{iu}-e^{iv}|\leq|u-v|$ gives
\[
                       |p(z+a)-p(z)|
                            \leq\frac{4M}{r-M}.
                                                               \tag{C.72.4}
\]
In particular, any plane homeomorphism $h$ for
which both $h(z)-z$ and $h^{-1}(z)-z$ are bounded
gives a closed-disk homeomorphism
\[
 \overline h(w)=
 \begin{cases}
       ph p^{-1}(w),&w\in\mathbb D,\\
       w,&w\in S^1 .
 \end{cases}                                                   \tag{C.72.5}
\]
For a compact continuous parameter family with a
common bound on both displacements, (C.72.4) gives
joint boundary continuity of the maps and inverses.
Interior continuity is just the actual conjugation
in (C.72.2). These observations prove continuity
in the full closed-disk homeomorphism topology.

\emph{The actual cyclic and translation actions.}
For $a\in\mathbb C$ put $L_a(z)=z+a$ and denote
its extension (C.72.5) by $\overline L_a$.
The equations
\[
 \overline L_a\overline L_b=\overline L_{a+b},
       \qquad\overline L_a^{-1}=\overline L_{-a}
                                                               \tag{C.72.6}
\]
hold first on the interior and then on the boundary.
The parameter map $a\mapsto\overline L_a$ is continuous
by (C.72.4), and its inverse on its image is
$h\mapsto p^{-1}(h(0))$. Its image is closed:
if $\overline L_{a_n}$ converges to a closed-disk
homeomorphism $h$, then $h(0)$ is interior and
$a_n\to p^{-1}(h(0))$. Thus this is a specified
closed translation subgroup, topologically $\mathbb R^2$,
acting transitively on $\mathbb D$ and fixing $S^1$
pointwise.

Set
\[
             T(x,y)=(x,y+1),\qquad
             T_*=\overline L_i,\qquad Q=\langle T_*\rangle.
                                                               \tag{C.72.7}
\]
Its interior image is closed and discrete: evaluating
the conjugate $T^m$ at the real origin gives $(0,m)$.
It is free and proper there, since compact real-plane
sets have bounded vertical coordinates and meet only
finitely many integer translates. These assertions
are transported by the actual $p,p^{-1}$.
Every element of $Q$ fixes the entire boundary.

Every forward orbit of $T_*$ has every point of
$S^1$ as a limit. Fix $z=x+iy$ and
$\zeta=e^{i\alpha}$. For large positive integers $k$
choose an integer $n_k$ nearest to
$\exp(2\pi k+\alpha-\pi/2)$. Then
\[
 \begin{split}
 \arg(x+i(y+n_k))&\longrightarrow\pi/2,\\
 \log(1+|x+i(y+n_k)|)
                -(2\pi k+\alpha-\pi/2)&\longrightarrow0,\\
 |p(x+i(y+n_k))|&\longrightarrow1.
 \end{split}                                                   \tag{C.72.8}
\]
The second limit follows since the integer rounding
error is at most $1/2$ and
$|x+i(y+n_k)|/n_k\to1$. Substitution in (C.72.2)
gives $p(T^{n_k}z)\to\zeta$.
Conversely $|T^n z|\to\infty$, so every accumulation
point is on the circle. This proves the stated
full-circle forward limit set for each original point.

\emph{A based loop with precisely this full local centralizer.}
Use the continuous functions
\[
 \varphi(y)=\operatorname{dist}(y,\mathbb Z),
                 \qquad\psi(x)=\max(0,1-|x|).
                                                               \tag{C.72.9}
\]
The first has zero set exactly $\mathbb Z$ and
period one; the second is nonzero and compactly
supported. For $0\leq u\leq1$ set
\[
 \begin{array}{ll}
 P_1(u)(x,y)=(x+u,y),&
 P_2(u)(x,y)=(x,y+u),\\
 P_3(u)(x,y)=(x+u\varphi(y),y),&
 P_4(u)(x,y)=(x,y+u\psi(x)).
 \end{array}                                                   \tag{C.72.10}
\]
Each is a homeomorphism, with inverse obtained by
replacing $u$ by $-u$ in the same formula. They
all commute with $T$. Their displacements and those
of their inverses have a common bound of one, so
their extensions $\overline P_j(u)$ in (C.72.5)
are continuous paths in
$C_{\mathcal D_\partial}(Q)$.

For $(j-1)/4\leq t\leq j/4$ define the triangular
amplitude
\[
 a_j(t)=
   \min\{\,8(t-(j-1)/4),\,2-8(t-(j-1)/4)\,\},
 \qquad f_t=\overline P_j(a_j(t)).
                                                               \tag{C.72.11}
\]
At every joining parameter the amplitude is zero,
so these are an actual continuous based loop.
Replacing $a_j(t)$ by $(1-r)a_j(t)$ gives its
explicit based contraction in the same relative
centralizer. In particular its ambient class is
zero, with no appeal to a different group.

Let $G=\langle f_t:0\leq t\leq1\rangle$.
It contains all real horizontal and vertical
translations in the $p$-coordinate: the intervals
$[0,1]$ in the first two paths generate the two
additive real groups, using products and inverses.
Thus $G$ is transitive on $\mathbb D$.
An interior homeomorphism centralizing $G$ must
commute with all these translations. Conjugation
by $p^{-1}$ and evaluation at zero then force it
to be a translation $L_{(a,b)}$.
Its commutation with $P_3(1)$ says
$\varphi(y+b)=\varphi(y)$ for every $y$.
At $y=0$ this gives $b\in\mathbb Z$.
Commutation with $P_4(1)$ says
$\psi(x+a)=\psi(x)$ for all $x$.
If $a\ne0$, iteration would give
$\psi(na)=\psi(0)=1$ for arbitrarily large $|na|$,
contradicting its compact support.
Hence $a=0$. Conversely these integer vertical
translations commute with every path (C.72.10).
We have proved the exact equality
\[
             C_{\operatorname{Homeo}(\mathbb D)}(G|_{\mathbb D})
                       =Q|_{\mathbb D}.
                                                               \tag{C.72.12}
\]
In particular the effective full local operator
group has dimension zero and is the named closed
discrete proper group, in precisely the topology
of the original comparison. The larger
$C_{\mathcal D_\partial}(Q)$ is also transitive,
since it contains the same translation subgroup.

\emph{The interior centralizer map that cannot extend.}
The real-plane map
\[
                        A(x,y)=(2x,y)
                                                               \tag{C.72.13}
\]
is a homeomorphism and commutes with $T$. Thus
$A_*=pAp^{-1}$ belongs to the target of (C.72.1).
Fix again $\zeta=e^{i\alpha}$ and, for sufficiently
large $k$, put
\[
 \begin{split}
 z_k&=(e^{2\pi k+\alpha}-1,0),\\
 z'_k&=(0,e^{2\pi k+\alpha-\pi/2}-1).
 \end{split}                                                   \tag{C.72.14}
\]
Both $p(z_k)$ and $p(z'_k)$ approach $\zeta$,
with the angular value exactly $\alpha$ modulo
$2\pi$. On the vertical sequence $A$ is the
identity, so $p(Az'_k)\to\zeta$. On the horizontal
sequence its image angle is
\[
 2\pi k+\alpha+
            \log(2-e^{-(2\pi k+\alpha)}).
                                                               \tag{C.72.15}
\]
Its radius tends to one, so
$p(Az_k)\to e^{i\log2}\zeta$.
These limits differ, since $0<\log2<1<2\pi$.
Thus $A_*$ has no continuous extension at $\zeta$.
As $\zeta$ was arbitrary, it has no continuous
extension at any boundary point. This proves the
failure of surjectivity in (C.72.1).

\emph{The specified interior Fuchsian comparison.}
Let $\mathbb H=\{u+iv:v>0\}$ and retain the actual
homeomorphisms
\[
 \begin{aligned}
 B(x,y)&=y+i e^x,& B^{-1}(u+iv)&=(\log v,u),\\
 C(\zeta)&=\frac{\zeta-i}{\zeta+i},&
 C^{-1}(w)&=i\frac{1+w}{1-w},\\
 j&=CBp^{-1}:\mathbb D\longrightarrow\mathbb D .
 \end{aligned}                                                   \tag{C.72.16}
\]
The displayed inverse formulas give
$j^{-1}=pB^{-1}C^{-1}$.
Direct substitution yields
$BTB^{-1}(\zeta)=\zeta+1$.
Hence $j$ conjugates $Q|_{\mathbb D}$ to the
parabolic Fuchsian group generated by
$C(\zeta\mapsto\zeta+1)C^{-1}$.
This interior conjugacy cannot extend to a
closed-disk homeomorphism: such an extension
would take the whole fixed boundary circle of
$T_*$ bijectively onto the single boundary
fixed point of that parabolic generator.
This is the same direct fixed-set equation as
(C.71.11), with both maps specified here.

For the obstructed map itself, the same
comparison gives
\[
                     BAB^{-1}(u+iv)=u+i v^2 .
                                                               \tag{C.72.17}
\]
This homeomorphism of $\mathbb H$ extends as
the identity on $\mathbb R\cup\{\infty\}$.
At finite real points this follows from
$v^2\to0$, and at infinity from
$|u+i v^2|\to\infty$ whenever
$|u+iv|\to\infty$. The inverse
$u+iv\mapsto u+i\sqrt v$ has the same property.
Consequently $jA_*j^{-1}$ is in the actual
boundary-fixed centralizer of the parabolic
closed-disk action, although $A_*$ is absent
from the original boundary-relative group.
Conjugation by $j$ is an isomorphism of the
interior centralizers, with inverse conjugation
by $j^{-1}$, but it does not carry the original
relative subgroup onto the parabolic one.
This proves the precise comparison failure
claimed in the statement.

All paths and their explicit contraction fix
the circle, so a specified closed-disk embedding
in a surface transports them and extends them
by the identity outside that disk. The inverse
families have the same extension. This places
the example in the original local-orbit setting
without changing its operator or loop labels.
The contraction in (C.72.11) continues to be
the actual contraction there.

The exact certificate
\path{code/part_b/c72_spiral_comparison_certificate.py}
checks the radial inverse and phase formulas,
the two displayed asymptotic phases symbolically,
the specified interior conjugacy and centralizer
map, and the periodic shear, inverse, and loop
identities with rational arithmetic.
The uniform boundary estimates, complete orbit
limit sets, full local-centralizer equality and
topological comparison failures are proved above.
\end{proof}
